\subsection{选择校准规模}

全池工作点得到\cref{tab:7-budget-full}中留出结果的支持；采样变体见\cref{fig:7-budget-risk}。\Cref{sec:6-2}的独立结果位模型提供规划尺度：$k=10$ 时，$4.74/(Sk)\approx10^{-3}$ 提示 $S$ 约为 500，而 128 的分辨率较粗。这一强模型针对参考不一致，不是 GT 损失，也不能将经验建议改写为预先指定的质量保证。

全池与留出批次使用同一扫描制度。十个单元采用距离截断；全池枚举已确认十二个实测单元均选到网格中最小接纳预算（\cref{sec:6-3}）。这一实测补充分组路径定理，但不将其扩展到任意距离并列。

\subsection{风险与所选预算性能}

共做 $m$ 次校准抽样，第 $r$ 次选择 $\widehat\alpha_r$，定义留出损失和越阈频率：
\begin{equation}
 L_r=\widehat r^*-\widehat r(\widehat\alpha_r),\qquad
 \widehat p_\delta=\frac{1}{m}\sum_{r=1}^{m}\mathbf1\{L_r>\delta\}.
 \label{eq:6-risk}
\end{equation}
平均损失、越阈频率和尾部大小回答不同问题。\Cref{tab:7-budget-quality,fig:7-budget-risk}报告归档摘要；补证逐抽样最大损失见\cref{tab:7-budget-tails}。

所选预算速率为 $\overline Q=m^{-1}\sum_r\widehat Q(\widehat\alpha_r)$，由已计时的固定预算查取。它既不同于中位预算的速率，也不同于校准后查询流的吞吐。质量与速度须一起解读（\cref{tab:7-budget-cost}）。

\subsection{复用成本}

对 $B_{\rm reuse}$ 个等大后续批次，模型为
\begin{equation}
 T_{\rm policy}(B_{\rm reuse})=T_{\rm cal}+B_{\rm reuse}T_{\rm chosen},\qquad
 T_{\rm fixed}(B_{\rm reuse})=B_{\rm reuse}T_f.
 \label{eq:6-total-cost}
\end{equation}
当 $T_{\rm chosen}<T_f$，回本批数为
\begin{equation}
 B_{\rm reuse}^*=\frac{T_{\rm cal}}{T_f-T_{\rm chosen}}.
 \label{eq:6-payback}
\end{equation}
严格超过 $B_{\rm reuse}^*$ 才有正节省；若选中批次不更快，则没有吞吐回本。对 $n$ 查询批，由 QPS 换算毫秒为 $T=1000n/Q$。选择变化时，$1/\mathbb E[Q]\leq\mathbb E[1/Q]$，因此使用 $1/\overline Q$ 偏乐观。实测调用边界见\cref{sec:a-budget-contract}；四个单元已用直接校准加查询流总成本验证这一模型，并同时报告质量（\cref{sec:7-stream-cost}）。\Cref{fig:7-budget-payback}给出当前实验估计。

\paragraph{一个复用成本算例。}
OpenAI-1536 在 128 探针时，全池校准选择 $\alpha=8$，请求 80 个保留候选，固定 100 则为 1,000。两者 recall 均为 0.9280，速率分别为 267,690 和 167,451 QPS。500 查询时，由此得到 $T_{\rm chosen}=1.868$\,ms、$T_f=2.986$\,ms；校准为 9.5\,ms，故 $B_{\rm reuse}^*=8.50$，随后九批有正模型节省。这些是归档预算的算术计算，不是直接计时的查询流成本。

若每 $H$ 批重新校准，会增加每批 $T_{\rm cal}/H$；漂移检测和自动重校准仍是后续工作。
